\begin{abstract}
Las Redes No Terrestres (NTN) representan un pilar fundamental para la cobertura global en 6G, pero enfrentan desafíos significativos en la gestión dinámica de recursos debido a la alta movilidad, variabilidad de canales y interferencias. Este artículo propone un marco basado en Inteligencia Artificial, específicamente Aprendizaje por Refuerzo Profundo (DRL) guiado por Modelos de IA Grandes (LAM), para la optimización automática de asignación de potencia, ancho de banda y beamforming en enlaces NTN. Mediante la formulación como un Proceso de Decisión de Markov (MDP) multi-agente, el sistema adapta recursos en tiempo real, maximizando el throughput sum-rate y el índice de Jain mientras minimiza la probabilidad de outage e interferencias co-canal. Se evalúa en escenarios realistas con constelaciones LEO, considerando propagación ITU-R y distribuciones heterogéneas de usuarios. Los resultados demuestran mejoras de hasta 40-64% en throughput y fairness frente a heurísticas y DRL baseline, con reducción significativa de interferencias. Se discuten implicaciones para integración TN-NTN y direcciones futuras en IA interpretable para NTN sostenibles.
\end{abstract}